On considère la suite $\left(u_{n}\right)$ définie par :
\[\left\{\begin{array}{l !{=} l}
u_{n + 1} & \dfrac{u_{n}}{1 + 2 u_{n}} \text { pour tout entier naturel } n 
\\ u_{0} & 1
\end{array}\right.\]

\begin{list}{\textbullet}{}
\item \textbf{Affirmation 1 :} \og{} $u_{4} = \dfrac{1}{9}$. \fg{}

$u_0=1$; 
$u_1=\dfrac{u_0}{1+2u_0} = \dfrac{1}{1+2\times 1}=\dfrac{1}{3}$;
$u_2=\dfrac{u_1}{1+2u_1} = \dfrac{\frac{1}{3}}{1+2\times \frac{1}{3}}
= \dfrac{\frac{1}{3}}{\frac{5}{3}} = \dfrac{1}{5}$;\\
$u_3=\dfrac{u_2}{1+2u_2} = \dfrac{\frac{1}{5}}{1+2\times \frac{1}{5}}
= \dfrac{\frac{1}{5}}{\frac{7}{5}} = \dfrac{1}{7}$;
$u_4=\dfrac{u_3}{1+2u_3} = \dfrac{\frac{1}{7}}{1+2\times \frac{1}{7}}
= \dfrac{\frac{1}{7}}{\frac{9}{7}} = \dfrac{1}{9}$

\hfill\textbf{Affirmation 1 vraie}

\item \textbf{Affirmation 2 :} \og{}Pour tout entier naturel $n$, $u_{n} = \dfrac{1}{2 n + 1}$.\fg{}

D'après la question précédente, la propriété \og $u_{n} = \dfrac{1}{2 n + 1}$ \fg{} est vérifiée pour $n$ entre 0 et 4. On va démontrer par récurrence qu'elle est vraie pour tout $n$.

\begin{list}{\textbullet}{}
\item \textbf{Initialisation}

Pour $n=0$, on a: $u_0=1$ et $\dfrac{1}{2\times 0 +1}=1$; donc la propriété est vérifiée.

\item On suppose la propriété vraie à un rang $n\geqslant 0$, c'est-à-dire que $u_n=\dfrac{1}{2n+1}$. C'est l'hypothèse de récurrence.

$u_{n+1}= \dfrac{u_n}{1+2u_n}
= \dfrac{\frac{1}{2n+1}}{1+2\frac{1}{2n+1}}
= \dfrac{\frac{1}{2n+1}}{\frac{2n+1+2}{2n+1}}
= \dfrac{1}{2n+1}\times \dfrac{2n+1}{2n+3}
=\dfrac{1}{2(n+1)+1}$

Donc la propriété est vraie au rang $n+1$.

\item \textbf{Conclusion}

La propriété est vraie au rang $0$, et elle est héréditaire pour tout $n\geqslant 0$, donc, d'après le principe de récurrence, elle est vraie pour tout $n$.
\end{list}

\hfill\textbf{Affirmation 2 vraie}

\item \textbf{Affirmation 3 :} \og{}La suite numérique $\left(u_{n}\right)$ est minorée par $10^{-10}$.\fg{}

$\ds\lim_{n\to +\infty} 2n+1=+\infty$ donc $\ds\lim_{n\to +\infty} \dfrac{1}{2n+1}=0$ et donc $\ds\lim_{n\to +\infty} u_n=0$.

Si la suite $(u_n)$ est minorée par $10^{-10}$ qui est un nombre strictement positif, elle ne peut pas avoir pour limite le nombre 0.

\hfill\textbf{Affirmation 3 fausse}

\end{list}
